\chapter{Tổng kết và Hướng phát triển}

\section{Tổng kết Chuỗi giá trị Khai phá (Pipeline Conclusion)}

Dự án "Dự báo năng suất cây trồng Khí tượng Nông nghiệp" đã khép kín trọn vẹn vòng đời khai phá tri thức nâng cao (End-to-End KDD). Toàn bộ nền tảng (Pipeline) được vận hành xuyên suốt không đứt đoạn: khởi điểm từ việc Sơ chế, lấp khuyết điểm dữ liệu (Notebook 01), tiến đến Khai phá chùm nguyên lý Sinh thái học bằng thuật toán FP-Growth/K-Means (Notebook 02), chuyển đổi đánh giá rủi ro thành bài toán Phân lớp đa luồng với Gradient Boosting (Notebook 03), và đạt đỉnh cao ở Mạng Nơ-ron Hồi quy chuỗi thời gian đan xen Không gian địa lý (Notebook 04 - LSTM).

Thành công cốt lõi về mặt kỹ thuật của dự án là việc "khai tử" triệt để hiện tượng chắp vá và lỗi rò rỉ kiến thức chéo (\textit{Data Leakage Illusion}) — một điểm mù thường xuyên đánh sập các mô hình học máy ứng dụng trong chuỗi thời gian. Toàn bộ kiến trúc phức tạp này được đóng gói minh bạch và tự động hóa dưới lớp điều phối \texttt{Papermill Auto-Execution}, đảm bảo tính tái lập (Reproducibility) tiêu chuẩn công nghiệp.

\section{Ý tưởng Thực thi Đắt giá (Actionable Insights)}

Từ góc nhìn Đánh giá và Phán quyết (Evaluate/Judgement), hệ thống trích xuất thành công 05 chỉ định chiến lược mang tính sống còn cho ngành Nông nghiệp:

\begin{enumerate}
    \item \textbf{Vạch trần Điểm Bão hòa Thuốc trừ sâu (Pesticides Plateau):}
    \\Luật kết hợp và Hồ sơ phân cụm (Tệp 1 vs Tệp 3) phát hiện một nghịch lý: Các vùng mang đặc tính "Nóng cực độ" (\textit{Climate Stress cao}) thường vung nạp siêu liều lượng thuốc sâu, nhưng sản lượng vẫn cắm đầu tụt dốc.
    \\$\Rightarrow$ \textit{Hành động (Góc độ Chính sách):} Chính phủ cần cắt ngay các gói trợ cấp hóa chất nông nghiệp ở vùng "Chảo lửa". Việc rải thêm hóa chất lúc này chỉ phá hoại vi sinh vật Đất, làm tăng chi phí vô ích chứ không thể đánh bật được thiên tai nhiệt độ. Cần chuyển hướng dòng vốn sang hệ thống tưới tiêu hoặc lai tạo giống.
    
    \item \textbf{Dịch chuyển Trọng tâm Cảnh báo (Trade-off Matrix Target):}
    \\Ma trận nhầm lẫn (Confusion Matrix) ở Notebook 03 cho thấy hệ thống có tỷ lệ \textit{False Positives} (Báo hỏng nhưng mùa không hỏng) dao động ở mức 20\%.
    \\$\Rightarrow$ \textit{Hành động (Doanh nghiệp/Chuỗi cung ứng):} Dưới góc độ quản trị rủi ro, người dùng phải chủ động chấp nhận 20\% báo động giả này để xả van đê điều trữ nước/tăng tồn kho, còn hơn đối mặt với 1\% \textit{False Negatives} (Hệ thống ru ngủ rằng mùa màng an toàn, nhưng thực tế đứt gãy chuỗi cung ứng dẫn đến khủng hoảng an ninh lương thực).
    
    \item \textbf{Cá nhân hóa Diện tích Chủng loài (Categorical Micro-segmentation):}
    \\Sức mạnh của mô hình phân lớp tăng vọt khi thay thế Label Encoding bằng ma trận phân mảnh \textit{One-Hot Category} (NB04). Điều này chứng minh toán học rằng: Quy luật sinh trưởng của Khoai tây vùng mưa lớn không thể chia sẻ chung không gian vector với Lúa nước.
    \\$\Rightarrow$ \textit{Hành động (Quy hoạch Agritech):} Yêu cầu lắp đặt hệ thống chip IoT đo lường vi khí hậu (Micro-climate) chuyên biệt cho từng đặc khu canh tác, tách biệt hoàn toàn module dự báo của từng loại cây trồng.
    
    \item \textbf{Khởi chạy Biến số Cảnh báo "Khắc nghiệt" Sinh học (Biological Stress Alerts):}
    \\Mô hình đã chứng minh các biến gốc (Lượng mưa, Nhiệt độ) rời rạc có sức mạnh dự báo thua xa các đặc trưng phái sinh giao thoa như \texttt{Climate\_Stress\_Index}.
    \\$\Rightarrow$ \textit{Hành động (Ứng dụng Di động):} Thay vì phát bản tin dự báo thời tiết cổ điển, các công ty Nông nghiệp Số (Agritech) cần tích hợp chỉ số \texttt{Climate\_Stress} vào App. Hệ thống sẽ tự động bắn thông báo (Push Notification) đến điện thoại nông dân ngay khi chênh lệch (Nhiệt - Mưa) vượt qua ranh giới chia cắt tối ưu (\textit{Splitting Threshold}) mà XGBoost đã tìm ra.
    
    \item \textbf{Định giá Tài sản Ảo giác và Cảnh giác Lừa đảo MAE (CV Illusion):}
    \\Biểu đồ \textit{CV Illusion Plot} đã lật tẩy sự dối trá của các ma trận đánh giá (Metrics) khi áp dụng Random Split cho dữ liệu có tính thứ tự thời gian.
    \\$\Rightarrow$ \textit{Hành động (Chủ đầu tư/Quỹ Khí hậu):} Loại bỏ lập tức các hồ sơ dự án AI Nông nghiệp sử dụng mô hình Cây (Trees) truyền thống nếu phát hiện họ xáo trộn (Shuffle) các năm học máy. Chỉ giải ngân chi trả cho các mô hình vượt qua được bài kiểm tra nghiệm thu \textit{Time-Series Split} khắc nghiệt (cắt đứt gãy quá khứ và tương lai).
\end{enumerate}

\section{Đề xuất Kiến trúc Cải tiến (Future Implementation Roadmap)}

Để đưa Pipeline từ một dự án Data Mining hàn lâm trở thành một sản phẩm MLOps thương mại (Production-ready), lộ trình nâng cấp kiến trúc vĩ mô bao gồm:
\begin{itemize}
    \item \textbf{Mở rộng Không gian Biến số Sinh - Hóa:} Tích hợp lượng Khí thải nhà kính ($CO_2$ equivalent) và độ màu mỡ phân bón đa lượng ($NPK$) thay vì chỉ giới hạn ở \texttt{Pesticides}, nhằm bao quát trọn vẹn chu trình Carbon trong nông nghiệp.
    \item \textbf{Tích hợp Cơ chế Tập trung Sâu (Deep Attention Mechanism):} Nâng cấp mạng LSTM hiện tại bằng kiến trúc Transformers/Attention. Thay vì đếm lùi tĩnh (\texttt{lookback = 3}), mạng lưới sẽ tự động tập trung trọng số (Attend) vào những năm có sự kiện thời tiết cực đoan (như El Nino/La Nina xảy ra cách đây 7 năm) có sự tương đồng sinh thái với năm hiện tại.
    \item \textbf{Tự động hóa Đám mây (Cloud Orchestration):} Triển khai hệ thống Cloud Cron-job để kích hoạt ngầm script \texttt{run\_papermill.py}. Hệ thống sẽ tự động thức giấc, cào dữ liệu mới nhất (Web Scraping) vào rạng sáng mùng 1 đầu năm từ API của tổ chức FAO, tự động re-train mô hình và gửi báo cáo rủi ro khẩn cấp qua Email cho các Bộ Nông nghiệp.
\end{itemize}